% Colored Text
\newcommand{\redtext}[1]{\textcolor{red}{#1}}
\newcommand{\orangetext}[1]{\textcolor{orange}{#1}}
\newcommand{\yellowtext}[1]{\textcolor{yellow}{#1}}
\newcommand{\greentext}[1]{\textcolor{green}{#1}}
\newcommand{\bluetext}[1]{\textcolor{blue}{#1}}
\newcommand{\purpletext}[1]{\textcolor{purple}{#1}}
\newcommand{\pinktext}[1]{\textcolor{pink}{#1}}
\newcommand{\violettext}[1]{\textcolor{violet}{#1}}

% Highlighted Text
\newcommand{\redbox}[1]{\colorbox{red}{#1}}
\newcommand{\orangebox}[1]{\colorbox{orange}{#1}}
\newcommand{\yellowbox}[1]{\colorbox{yellow}{#1}}
\newcommand{\greenbox}[1]{\colorbox{green}{#1}}
\newcommand{\bluebox}[1]{\colorbox{blue}{\textcolor{white}{#1}}}
\newcommand{\purplebox}[1]{\colorbox{purple}{\textcolor{white}{#1}}}
\newcommand{\pinkbox}[1]{\colorbox{pink}{#1}}
\newcommand{\violetbox}[1]{\colorbox{violet}{#1}}

%%% Center Text %%%
\newcommand{\centertext}[1]
{%
    \begin{center}
        #1
    \end{center}
}

% ===== Tables =====

% Make a table with a header
% #1 is |c|c|c|
% #2 is H1 & H2 & H3 \\
% Body is B1 & B2 & B3 \\
\newenvironment{tablewithheader}[2]
{% begin
    \begin{center}

        \begin{tabular}{#1}

            % Header
            \hline
            #2 \\ [0.5ex]
            \hline

            % Body
            % =====
            \hline

}
{% end

            \hline
            % =====

        \end{tabular}
    \end{center}
}

% Make a table with a header aligned to the left
% #1 is |c|c|c|
% #2 is H1 & H2 & H3 \\
% Body is B1 & B2 & B3 \\

% EXAMPLE
% \begin{tablewithheader}
% { |c|c|c| }
% { Col A & Col B & Col C }

%     1 & 2 & 3 \\
%     a & b & c \\ 
%     x & y & z \\

% \end{tablewithheader}

\newenvironment{tablewithheaderleft}[2]
{% begin

    \begin{tabular}{#1}

        % Header
        \hline
        #2 \\ [0.5ex]
        \hline

        % Body
        % =====
        \hline

}
{% end

        \hline
        % =====

    \end{tabular}
}

% Make a table without a header
% #1 is |c|c|c|
% Body is B1 & B2 & B3 \\
\newenvironment{tablewithoutheader}[1]
{% begin
    \begin{center}

        \begin{tabular}{#1}

            % Body
            % =====
            \hline

}
{% end

            \hline
            % =====

        \end{tabular}
    \end{center}
}

% FFFF00 (yellow)
% FFFF9E (yellow)
% When using |c|c|c| replace "c" with "s" for a highlighted column
\newcolumntype{s}{>{\columncolor[HTML]{FFFF9E}} c}

% ===== Lists =====

% Numbered List
% 
% \begin{enumerate}
%     \item ...
% \end{enumerate}

% Bullet-point List
% 
% \begin{itemize}
%     \item ...
% \end{itemize}

% Alphabet List
\newlist{alphabet}{enumerate}{1}
\setlist[alphabet]{label=(\alph*)}

% Roman Numeral List
\newlist{romanlist}{enumerate}{1}
\setlist[romanlist]{label=(\roman*)}

% ===== Math =====

% Multi-lined Math
% 
%   - Everything aligns at &
% 
% \begin{align*}
%     a + b &= c \\
%     &= 2 + 2 \tag{ This is a tag! } \\
%     &= 4 \\
% \end{align*}

% ========== Set Notation ==========

% Italics Shorthand
\newcommand{\mathset}[1]
{
    \ensuremath{\mathit{#1}}
}

% Brackets Shorthand
\newcommand{\mathinset}[1]
{
    \ensuremath{\{ #1 \}}
}

% Define colors for set notation
\newcommand{\definesetcolors}{%

    % color of set of Natural numbers
    \colorlet{colN}{Red!60}
    % color of set of Whole Numbers
    \colorlet{colW}{OrangeRed!60}
    % color of set of Integers
    \colorlet{colZ}{Orange!60}
    % color of set of Rational numbers
    \colorlet{colQ}{Gold!60}
    % color of set of Real numbers
    \colorlet{colR}{ForestGreen!50}
    % color of set of Complex numbers
    \colorlet{colC}{CornflowerBlue!80}
    % color of set of Hypercomplex numbers
    \colorlet{colH}{BlueViolet!50}
    % color of set of Transcendental numbers
    % \colorlet{colT}{Green!30}

    % center coordinate of Natural numbers
    \coordinate (oN) at (0.0, 0.0);
    % center coordinate of Whole Numbers
    \coordinate (oW) at (0.5, 0.5);
    % center coordinate of Integers
    \coordinate (oZ) at (1.0, 1.0);
    % center coordinate of Rational numbers
    \coordinate (oQ) at (1.5, 1.5);
    % center coordinate of Real numbers
    \coordinate (oR) at (2.0, 2.0);
    % center coordinate of Complex numbers
    \coordinate (oC) at (2.5, 2.5);

    % set of Natural numbers
    \def\setN{(oN) ellipse (2.5 and 1.0)}
    % set of Whole numbers
    \def\setW{(oW) ellipse (3.5 and 2.0)}
    % set of Integers numbers
    \def\setZ{(oZ) ellipse (4.5 and 3.0)}
    % set of Rational numbers
    \def\setQ{(oQ) ellipse (5.5 and 4.0)}
    % set of Real numbers
    \def\setR{(oR) ellipse (6.5 and 5.0)}
    % set of Complex numbers
    \def\setC{(oC) ellipse (7.5 and 6.0)}
}

% ========== Discrete Specific ==========

\newcommand{\T}{\greentext{T}}
\newcommand{\F}{\redtext{F}}

\newcommand{\TRUE}{\greentext{TRUE}}
\newcommand{\FALSE}{\redtext{FALSE}}